\RequirePackage[hyphens]{url}
\documentclass[11pt,titlepage,dvipsnames,table,xcdraw]{article}
\usepackage{strath-dissertation}
\usepackage{setspace}

\usepackage{babel}
\usepackage{lipsum} % For demonstration only.

% packages specific to this document
\usepackage{pdfpages}
\usepackage{listings}
\usepackage{adjustbox}

%--------------------------------
% Degree settings
%--------------------------------
\newcommand{\degreename}{Bachelor of Science}
\newcommand{\coursename}{CS408: Individual Project}
\newcommand{\deptname}{Computer and Information Sciences}

\begin{document}

%--------------------------------
% Title page.
%--------------------------------
\begin{titlepage}
\vspace*{5mm}
\titletext{LinTeX:\\Promoting Accessibility in\\LaTeX Documents}
\vspace{10mm}
\authortext{Mohammad Danyal}
\vspace{15mm}
\centering{\Large{\textbf{Progress Report}}}
\vspace{40mm}
\centredcrest
\vspace{5mm}
\depttext{\deptname}
\vspace{5mm}
\datetext{November, 2025}
\end{titlepage}

%--------------------------------
% Front matter numbering.
%--------------------------------
\pagenumbering{roman}
\setcounter{page}{2}

%--------------------------------
% Setting counter depth.
%--------------------------------
\setcounter{secnumdepth}{3}
\setcounter{tocdepth}{2}

%--------------------------------
% The front matter.
%--------------------------------
\setstretch{1.2} % 1.5 line spacing. 




%--------------------------------
% Tables of contents.
%--------------------------------
\setstretch{1.0} % 1.0 line spacing.  
\clearpage
\tableofcontents
%\listoffigures   % Uncomment for a list of figures.
%\listoftables    % Uncomment for a list of tables.
\clearpage

%--------------------------------
% The body of the document.
%--------------------------------
\pagenumbering{arabic} % Page numbering for rest of document.
\setstretch{1.5} % 1.5 line spacing. 


\textbf{The page count of 12 is exclusive of bibliography and appendices.}
\section{Aims and Objectives}
% Aims and Objectives:
% + Say what is useful and difficult about your project
% + Use a mix of requirements engineering, hypotheses, and research
%     questions to set you aims — these have an overlap with the "Project
%     Specification" section.

% approximately one page

There is a lack of accessibility with LaTeX, largely due to usage of outdated PDFLaTeX as opposed to LuaLaTeX. Usage of this outdated TEX engine necessitates workarounds for desired features and accessibility in resulting PDF documents.

\begin{quote}
What can be done to make LaTeX documents more accessible?
\end{quote}
LinTeX is a developing tool dedicated to combat this problem by focusing on assisting screen readers' abilities to read the PDF documents produced from LaTeX documents. It aims to cover this area through suggesting structured tagging throughout a document, suggesting components where it may be lacking. Additionally, it will request users to input appropriate alternative text ('alt text') for non-text components of the document, including figures, images, and tables. Further focus may be placed on features such as ensuring appropriate contrast in any colours in the document, appropriate Sans Serif fonts, and flexible document design, largely depending on what development time allows for. 

Additionally, as a tool focused on promoting accessibility, it too should be accessible. Therefore, the tool will ensure strict adherence to the Web Content Accessibility Guidelines 2 (WCAG 2) \cite{wcag} requirements with a focus placed on keeping it a simple and intuitive tool for users to make use of.

\begin{quote}
Why should I use LinTeX?
\end{quote}
While the tool is not limited strictly to the field, it is aimed to assist the many academics that find difficulties in accessing reports and academic literatures. By utilising this tool, you can ensure a wider demographic of users can view your documents, without necessitating an especially large amount of research in accessibility; it has been done for you.
\clearpage

\section{Summary of Related Work}
% Summary of Related Work:
% This can be covered by the competitive evaluation of shopping lists of
% criteria for the candidates, plus identifying "things" that have
% influenced your decisions / design. For example if you are making a
% game, then you would review other games that are influencing you but
% limited to those aspects (so not writing a magazine / website review of
% a game).

% You may summarise one or more former project student's work if it's
% relevant (from the Edinburgh lists I posted).

% All of you will want to cite guidelines for designing the user
% interface; with additional cites for accessibility-related matters; and
% further cites for designing your user study.

% approximately 5 pages

\subsection{Programming Language Comparisons}
There were many considerations when deciding what programming language to write LinTeX in. Languages chosen were based on those that I had prior experience in, to ensure minimal hinderance from the language itself to the development of the project. An important factor is ensuring LinTeX's maintainability in the future. Therefore, only the languages which have actively maintained Parsing Expression Grammar (PEG) libraries will be considered, as can be seen in Table 1.1. The most recent update times provided in Table 1.1 are in reference to 24\textsuperscript{th} November 2025.
\begin{center}
Table 1.1: Languages and Related PEG Frameworks
\begin{adjustbox}{width=\textwidth}
\begin{tabular}{ c | c | c}
Language & PEG Library Link & PEG Library is in Development \\
\hline
Python & https://pythonhosted.org/pyPEG2/ & No, last updated in October 2015 \\
Python & https://github.com/erikrose/parsimonious & Yes, last updated in November 2025 \\
JavaScript & https://github.com/pegjs/pegjs & No, last updated in November 2019 \\
JavaScript & htps://peggyjs.org/ & Yes, last updated in September 2025 \\
Go & https://github.com/pointlander/peg & Yes, last updated in November 2025 \\
Go & https://github.com/mna/pigeon & Yes, last updated in November 2025 \\
C & N/A & N/A \\
Java & https://github.com/sirthias/parboiled & Yes, last updated in November 2025 \\
C\# & https://code.google.com/archive/p/peg-sharp/ & No, last updated in May 2011 \\
C\# & https://github.com/otac0n/Pegasus & No, last updated in December 2023 \\ 
C++ & https://github.com/yhirose/cpp-peglib & Yes, last updated in July 2025 \\
PHP & https://github.com/hafriedlander/php-peg & No, last updated in May 2014 \\
TypeScript & https://github.com/metadevpro/ts-pegjs & No, last updated in November 2024 \\
TypeScript & https://github.com/EoinDavey/tsPEG & Yes, last updated in November 2025 \\
Dart & https://github.com/darmie/dart-peg & No, last updated in July 2015 \\
\end{tabular}
\end{adjustbox}
\end{center}
Based on the current criteria, we can rule out C, C\#, PHP, and Dart.

As many languages remain, there is a need for a stricter criteria: suitability to be utilised as a web application as well as compatibility with the University of Strathclyde's development web server (devweb) system.
\begin{center}
Table 1.2: Languages and Web Development Support
\begin{adjustbox}{width=\textwidth}
\begin{tabular}{ c | c | c}
Language & Web Development Support & DevWeb Support \\
\hline
Python & Yes & Yes \\
JavaScript & Yes & Yes \\
Go & Yes & No \\
C & No, except for rare exceptions & No \\
Java & Yes, but not generally used & Yes \\
C\# & Yes, but not generally used & No \\ 
C++ & Yes, but not generally used & No \\
PHP & Yes & Yes \\
TypeScript & Yes & Yes \\
Dart & Yes & Yes \\
\end{tabular}
\end{adjustbox}
\end{center}
It should be noted that while Go, C, C\#, and C++ are not directly compatible with devweb, as referenced in Table 1.2, they can be supported through use of Common Gateway Interface (CGI) scripts, and consequently should not be fully disregarded for this alone. However, based on their support for web development, C, Java, C\#, and C++ will have a lower priority for being a chosen language.

This leaves us with a focus placed on JavaScript, Go, and TypeScript. However, JavaScript will be removed from consideration as a strictly typed language is preferred for our context. Therefore, leaving Go and TypeScript as the main considerations for this project. Finally, as I have more experience and a stronger preference towards Go as opposed to TypeScript, it was the chosen programming language for the project.

\subsection{Library Comparison}
Go has two libraries that are used in relation to PEG systems: Pigeon by Martin Angers (mna), and peg by Andrew Snodgrass (pointlander). Direct comparisons with core components of both have been made in Table 1.3.
\begin{center}
Table 1.3: Comparison Between PEG and Pigeon
\begin{adjustbox}{width=\textwidth}
\begin{tabular}{ c | p{100mm} p{100mm}}
& Pigeon & Peg \\
\hline
Installation & Requires running a \lstinline[language=Go]!go install! command and setting up environment variables & Requires supporting bash scripting (through Window's Subsystem for Linux (WSL) for Windows users), as well as running a \lstinline[language=Go]!go install! command, cloning the repository, and running \lstinline[language=Go]!go generate \&\& go build! \\
Documentation & Contains a Godoc page, wiki page, an examples subdirectory, and a detailed README.md providing what PEG features are supported, command line interface usage, installation, and examples & Contains a detailed README.md and a PEG file syntax file containing what PEG features are supported, installation steps, and example projects links \\
Generated Code & Generated Go code can be used to configure settings on the parser, view statistics or error information, debug, enable or disable memoization, and parse readers, files, and strings. Additionally, the generated code is well documented, with several comments surrounding usage of public functions & Generated Go code can be used to run the PEG script on a provided string, debug and output the Abstract Syntax Tree (AST), output errors in parsing, and enable or disable memoization \\
Examples & Provides three direct examples: detecting indentation in files, parsing JSON files, and a calculator example. Examples were kept short (in regards to PEG scripts) with the calculator and indentation examples showcasing usage via a main function & Provides two project links for examples on PEG usage: a natural date/time parser and a scientific names parser. Both projects have large PEG scripts, with most lines being simple to read. The natural date/time project was intuitive to use and understand program flow \\
\end{tabular}
\end{adjustbox}
\end{center}
While examples are relatively similar, the simplicity of installation for Pigeon as opposed to Peg, combined with the depths of documentation both in generated code and outwith it, makes Pigeon a more ideal option. Additionally, its greater number of features gives LinTeX more to make use of. Therefore, Pigeon was the chosen library for the project.

\subsection{Design Focus}
LinTeX will take strong inspiration from Integrated Development Environments (IDEs) with regards to its User Interface/User Experience (UI/UX) design. This is because LaTeX in structure, is similar to a programming language, even being Turing complete \cite{latexturingcomplete}. Hence, an environment suitable to programming is suitable to LaTeX, with some special considerations put aside for its quirks. For example, depending on the user's screen size and preferences, they may wish to have the output of the LaTeX documents taking up half of the screen, a separate browser window or monitor, or out of the way until compilation. 

% if you're reading this, you owe me maryland cookies
% \subsection{TABLE FOR REFERENCE}
% \begin{center}
% Table 1.1: Choosing a Programming Language
% \begin{adjustbox}{width=\textwidth}
% \begin{tabular}{ c | c | c | c | c | c | c}
% Language & Experience Level & Web Development Support & PEG Framework Exists & PEG Framework in Development & Personal Preference Prioritization & DevWeb Support \\
% \hline
% Python & Medium & Yes & https://pythonhosted.org/pyPEG2/ \newline https://github.com/erikrose/parsimonious & No, PyPEG2 was last updated in October 2015 \newline Yes, parsimonious was last updated in November 2025 & 5 & Yes \\
% JavaScript & High & Yes & htps://peggyjs.org/ & Yes, last updated in September 2025 & 4 & Yes \\
% Go & Medium in Go, High in \href{https://pkg.go.dev/text/template}{Text/Template Scripting Language} & https://github.com/pointlander/peg & Yes, last updated in November 2025 & 1 & Yes, as CGI scripts \\
% C & Low & No, except for rare exceptions & N/A & N/A & 8 & No \\
% Java & Medium & Yes, but not generally used & https://github.com/sirthias/parboiled & Yes, last updated in November 2025 & 7 & Yes \\
% C\# & High & Yes, but doesn’t seem great & https://github.com/otac0n/Pegasus & No, last updated in December 2023 & 2 & No \\ 
% C++ & Low & Yes, but doesn’t seem great & https://github.com/yhirose/cpp-peglib & Yes, last updated in July 2025  & 9 & No \\
% PHP & Medium & Yes & https://github.com/hafriedlander/php-peg & No, last updated May 2014 & 10 & Yes \\
% TypeScript & Medium & Yes & https://github.com/EoinDavey/tsPEG  & Yes, last updated in November 2025 & 3 & Yes \\
% Dart & Low & Yes & https://github.com/darmie/dart-peg  & No, last updated July 2015 & 6 & Yes \\
% \end{tabular}
% \end{adjustbox}
% \end{center}
\clearpage

\section{Project Specification}
% approximately 2 pages

Modern LaTeX documents utilise PDFLaTeX as opposed to LuaLaTeX. Thus, they often require workarounds to implement features and lack accessibility support. LinTeX is a new tool under development to combat this concern, which aims to improve LaTeX documents by avoiding workarounds and suggesting modern LaTeX usage as an alternative, with emphasis placed on accessibility that support screen readers. 

Functional requirements include:
\begin{itemize}
\item LinTeX must allow one LaTeX file to be uploaded.
\item LinTeX should allow multiple LaTeX files to be uploaded.
\item LinTeX should allow non-LaTeX files (such as PDFs) to be uploaded.
\item LinTeX must allow the files to be exported as a ZIP file once the user has finished using the system.
\item LinTeX should allow the finished document to be previewed.
\item LinTeX must detect if tagging exists or not.
\item LinTeX must provide a fix by inserting tagging.
\item LinTeX must detect when alt text is missing for images.
\item LinTeX must detect when alt text is missing for figures.
\item LinTeX must detect when alt text is missing for tables.
\item LinTeX must provide a fix by inserting alt text for images.
\item LinTeX must provide a fix by inserting alt text for figures.
\item LinTeX must provide a fix by inserting alt text for tables.
\item LinTeX should support batch error fixes for errors of the same type.
\item LinTeX should reference further reading for any accessibility recommendations.
\item The website should have all non-text content contain an alt text (WCAG 2.2: Guideline 1.1) \cite{wcag}.
\item LinTeX should not be vulnerable to any of OWASP's 2024 top 10 vulnerabilities.
\item LinTeX may be usable without a keyboard.
\item LinTeX should provide simple error messages.
\item LinTeX should provide simple suggestion messages.
\end{itemize}

Non-functional requirements include:
\begin{itemize}
\item The website should load within 2 seconds.
\item The website should support a variety of standard desktop screen sizes.
\item The website should support a variety of browsers.
\item The website should be presentable in a simpler format if required (WCAG 2.2: Guideline 1.3) \cite{wcag}.
\item The website should allow for content to be distinguishable including supporting resizing content and minimal colour usage (WCAG 2.2: Guideline 1.4) \cite{wcag}.
\item The website may be operable through the keyboard alone (WCAG 2.2: Guideline 2.1) \cite{wcag}.
\item The website should include navigation features (WCAG 2.2: Guideline 2.4) \cite{wcag}.
\item The website should be predictable (WCAG 2.2: Guideline 3.2) \cite{wcag}.
\item The website should follow the law of locality.
\item LinTeX should provide a PDF document for several files in under a minute.
\item LinTeX should create a ZIP of the documents in under a minute.
\item LinTeX should have thorough documentation describing its usage and code.
\item LinTeX should be able to handle several files with hundreds of lines each.
\item LinTeX should be able to handle thousands of users using the system simultaneously.
\end{itemize}
\clearpage

\section{Project Plan including summary of progress}
% Project Plan including summary of progress:
% A place to show your end-to-end development, and perhaps something about
% the user study.

% For the plan part of the report, I do want to see some indication of
% dates and durations, eg you will spend the second two weeks of January
% doing a user test.

% ideally includes some prototype development (approximately 2 pages)
\subsection{Future Plans}
Focus will be shifted to completing the ethics application from 1\textsuperscript{st} December 2025 to 14\textsuperscript{th} December 2025. Once completed, it will be aimed to complete the user study during the second half of January, to give appropriate time to make any necessary adjustments to the ethics application.

When the ethics application is complete, focus will be shifted to completing the prototype and relevant research for it. Predictably, this would begin on 15\textsuperscript{th} December 2025 or sooner. Ensuring basic functionality, with a focus on backend, will occur in the initial stages, likely requiring a week or two of dedicated efforts. Afterwards, Figma wireframes will be produced for the UI/UX to ensure it meets accessibility requirements which may take a week to complete. After suitable Figma diagrams have been created, implementation of the wireframes will take place which may take up to two weeks, due to a lack of experience with Tailwind CSS.

\subsection{End-to-End Prototype}
A basic webpage was created in Go utilising the Go package template (html/template), as showcased in Figure 4.1. This website was created following a tutorial on the Go webpage \cite{gowebtutorial}. Additionally, it incorporated basic styling to confirm CSS could be ran. 
\begin{center}
Figure 4.1: Basic Website\\
\includegraphics[width=0.5\textwidth]{images/basic_website.png}
\end{center}
After confirming a basic webpage could be created, the next objective was confirming a PEG script could be ran to generate Go code. This was confirmed utilising a basic calculator example PEG script provided in Pigeon documentation \cite{pigeon}, which ran without issues as confirmed in Figure 4.2.
\begin{center}
Figure 4.2: Generating Go Code\\
\includegraphics[width=0.8\textwidth]{images/generate_go_code.png}
\end{center}
It was important to confirm the generated Go code was usable, which was confirmed by utilising a simple function to receive input from the user (the calculation to be performed), and outputed the correct amount by making use of test.go.
\begin{center}
Figure 4.3: Running Basic PEG Script\\
\includegraphics[width=0.7\textwidth]{images/testing_peg_script.png}
\end{center}
Immediate next steps would be creating a basic PEG script which prints an output to console when specific LaTeX tags are detected. Additionally, it will be important to ensure it can be connected to the website to be ran. Therefore, a further step includes receiving user input from the website, running generated Go code from the PEG script on it, and outputing relevant information to the website. This would work as a minimum proof of concept, as it would confirm we can identify a missing LaTeX tag, for example missing alt text associated with an image.

\subsection{User Study}
A prototype Consent Form and Participant Information Sheet have been created, as can be viewed in the appendix. Next steps include finalising details for the ethics application as well as receiving further reviews.
\clearpage

\section{Approach to Development}
% Brief summary of proposals for Development Methodology, Design, Implementation, Testing
% and Evaluation including proposed technologies to be used (approximately 2 pages).
% + Methodology:
%    - more likely this is a method or technique rather than a methodology.
% + Design
%    - your overview diagram, screenshots, and indicating how you meet
%      guidance on interface design.
% + Implementation and Technology ideas
%    - refer back to the competitive evaluation from earlier, perhaps here
%      summarising the winning combination of parts.
% + Testing and Evaluation plans
%    - you can use Net Promoter Score and statistical significance testing
%      on results (see Trista Yang's TriTeX MSc for all of that);
%    - you can say more about your user study;
%    - you can talk about data design, level design (for games), regression
%      testing, and any other Software Engineering techniques that are
%      relevant — remember: despite all the design work and creativity,
%      this is still a Computer Science project and needs to be easily
%      identified as a Computer Science project.

\subsection{Methodology}
Test-Driven Development (TDD) is the chosen development methodology for this project. The reason for this is that there are a large number of features in LaTeX and involving testing at every stage ensures that each component can be handled in a methodological manner. Additionally, this is an opportunity to explore a development methodology that I have not had the opportunity to utilise beforehand, so it is an opportunity to learn.

TDD is a development methodology which often works in sprints similar to an Agile Development Methodology. However, it differs in that the test cases are written first and the code is then adapted to pass the test cases \cite{tdd}. Therefore, it is imperative that the test cases written for LinTeX are thorough.

\subsection{Design}
The initial overview diagram, as showcased in Figure 5.1, was used to showcase a document lacking accessibility becoming more accessible after using LinTeX.
\begin{center}
Figure 5.1: Overview Diagram Prototype\\
\includegraphics[width=\textwidth]{images/overview_diagram_draft_1.png}
\end{center}
However, the prototype required the images within it to showcase the accessibility in a strong manner, which could be difficult to do with them alone. Therefore, the overview diagram has become more focused on a panel-based design (i.e. a 'comic') utilising characters from the UK government's disability and impairment user profiles \cite{govuserprofiles}. There are several characters involved: Ashleigh, a partially sighted screenreader user; Christopher, a user with rheumatoid arthritis; and LinTeX, a printer-like character which takes LaTeX documents and makes them more accessible.

Ashleigh and Christopher will initially be having a conversation about difficulty with reading academic documents, with LinTeX easedropping on their conversation. LinTeX will then offer a solution by taking their document and putting it through his printer-like services to create a more accessible end document. Finally, Ashleigh and Christopher will comment on how it has made the document easier to read. For example, mentioning better accuracy with supporting screen readers.

\subsection{Implementation and Technology Ideas}
Go is the programming language of choice. A PEG library has been chosen as core to the project, with Pigeon being the chosen library to implement it \cite{pigeon}. Go package template (html/template) will be used to implement the webpages \cite{gotemplate}, with Tailwind CSS used to create the styling. CGI scripts on devweb or an nginx web server will be used to run the code, with further considerations being necessary. Go package testing will be used to implement test cases \cite{gotesting}, with generated test cases ('fuzzing') used where possible.

\subsection{Testing and Evaluation Plans}
To evaluate the accessibility of the LaTeX document, and usability of LinTeX, a user study will be conducted. It will have several steps included: participants will be asked to undergo a think-aloud session where they will use LinTeX, as well as answer some questions about their experience after completing the think-aloud session. A net promoter score could be generated based on the questions asked. 
\clearpage

\clearpage

%--------------------------------
% Bibliography
%--------------------------------
\clearpage
\setstretch{1.0}  % 1.0 line spacing. 
\bibliographystyle{apalike}
\bibliography{bibliography}
% add references to the file called bibliography.bib
\clearpage

%--------------------------------
% Appendix
%--------------------------------
\appendix
\includepdf[pages=-]{images/participant_info_sheet.pdf}
\includepdf[pages=-]{images/consent_form.pdf}


\end{document}